% =====================================================================
% Section 4: Drag Models
% Parent jsr_paper.tex already declares \section{Drag Models} and \input's this file.
% Start at \subsection. Equation numbering is automatic; do not hard-code \tag.
% =====================================================================

The total axial-force coefficient is assembled from independently validated
contributions, each evaluated at the locally corrected post-shock state supplied
by the shock-geometry pre-pass of Sec.~II:
%
\begin{equation}
C_D = C_{D,\,\mathrm{friction}} + C_{D,\,\mathrm{pressure}}
    + C_{D,\,\mathrm{base}} + C_{D,\,\mathrm{override}}
    + C_N \sin\alpha .
\label{eq:cd-total}
\end{equation}
%
Each submodel below states its regime of applicability, governing equation, and
the published benchmark against which it was assessed. Table~\ref{tab:drag-inventory}
summarizes the inventory.

\subsection{Nose and Body Wave Drag}

For conical noses the wave-drag coefficient equals the Taylor--Maccoll surface
pressure coefficient---the exact inviscid result for steady conical flow at zero
incidence (Sec.~III)---and serves as the reference for all shape-correction
methods.

For tangent and secant ogives, parabolic, and shock-attached power-law noses, a
shock-expansion strip integrator marches $N_{\mathrm{strip}} = 100$ conical
frustum strips from tip to base. Each strip applies a Prandtl--Meyer expansion or
an oblique shock according to the local turning angle; pressure and temperature
ratios accumulate multiplicatively. The pressure-drag integral is
%
\begin{equation}
C_{d,\,\mathrm{wave}} = \frac{2}{R_{\mathrm{aft}}^2 - R_{\mathrm{fore}}^2}
   \sum_{i=1}^{N_{\mathrm{strip}}} C_{p,i}\,r_{\mathrm{mid},i}\,\Delta r_i ,
\label{eq:strip-wave}
\end{equation}
%
summing only windward strips, with tip initial conditions seeded from
Taylor--Maccoll at the local tip half-angle.

For Haack-series, parabolic, and selected power-law noses, the Dahlem--Buck
correction~\cite{dahlembuck1966} scales the equivalent-cone result by an empirical
shape factor and a fineness correction,
%
\begin{equation}
C_{d,\,\mathrm{wave}} = C_{d,\,\mathrm{cone}}(M, \theta_{\mathrm{equiv}})\,
   K_{\mathrm{shape}}\,\left(\frac{3}{f}\right)^{1.6} ,
\label{eq:dahlem-buck}
\end{equation}
%
with $K_{\mathrm{shape}} = 1.00$ for cones, $0.85$ for Lighthill--von~K\'arm\'an
(L-V) ogives, $0.88$ for parabolics, $0.90$--$0.95$ for power-law families, and
$0.60$ for the Haack series. A smoothstep over $M = 1.3$--$1.5$ fades the legacy
NASA~TR~R-100~\cite{nasa_trr100} tabulated drag into the analytical result above
$M = 1.5$. Below the drag-divergence Mach
$M_{\mathrm{dd}} = 0.95 - 0.15\sin^{0.4}\theta_{\mathrm{tip}}$ wave drag is zero;
above $M_{\mathrm{dd}}$ a $C^1$ cubic Hermite connects zero drag to the first
analytical point. Above $M = 5$ the nose pressure crosses over to Modified
Newtonian theory (Sec.~IV.F). The combined model was validated against
NACA~RM~A52H28~\cite{a52h28} wind-tunnel pressure measurements for five
fineness-3 nose shapes (cone, quarter-power, three-quarter-power, Haack, and L-V
ogive) at $M = 1.5$--$3.0$, giving an aggregate mean absolute error of $0.029$ in
$C_d$ against an acceptance gate of $0.035$ (Fig.~\ref{fig:a52h28}).

\begin{figure}[t]
\centering
\includegraphics[width=\columnwidth]{naca_rm_a52h28_validation}
\caption{Nose wave drag versus NACA~RM~A52H28~\cite{a52h28} for cone,
quarter-power, three-quarter-power, Haack, and L-V ogive noses at fineness~3,
$M = 1.5$--$3.0$. Aggregate mean absolute error $= 0.029$ in $C_d$.}
\label{fig:a52h28}
\end{figure}

\subsection{Fin Wave Drag (DATCOM 4.1.5.1)}

The fin wave-drag model replaces the legacy $\cos^2\Lambda_{LE}$ Ackeret
approximation with the full DATCOM~4.1.5.1 method~\cite{datcom1978}, branching on
whether the fin leading edge is subsonic or supersonic relative to the local
post-shock Mach number:
%
\begin{equation}
C_{d,\,\mathrm{wave}} = \begin{cases}
\dfrac{K}{\beta}\left(\dfrac{t}{c}\right)^2,
   & \beta\cot\Lambda_{LE} \ge 1 \quad (\text{supersonic LE}),\\[2.0ex]
K\cot\Lambda_{LE}\left(\dfrac{t}{c}\right)^2,
   & \beta\cot\Lambda_{LE} < 1 \quad (\text{subsonic LE}),
\end{cases}
\label{eq:datcom-fin}
\end{equation}
%
where $\beta = \sqrt{M^2 - 1}$, the section shape factor $K = 4.0$ for
double-wedge (hexagonal) sections and $16/3$ for biconvex or rounded sections per
DATCOM Table~4.1.5.1-A, and $t/c$ is the streamwise thickness-to-chord ratio. A
cubic Hermite blend over $M = 0.9$--$1.2$ fades zero wave drag into the full
DATCOM result. The implementation reproduces the closed-form Ackeret formula on
15~analytical cases to $0.00\%$ error (code verification) and was validated
against free-flight measurements of a 60-deg delta wing from
NACA~TN~3650~\cite{ulmann1956}, yielding a mean absolute percentage error of
$21.0\%$ across 12~points over $M = 1.1$--$2.5$ (Fig.~\ref{fig:fin-wave}). The
residual reflects the difficulty of predicting wave drag near the
subsonic-to-supersonic leading-edge transition.

\begin{figure}[t]
\centering
\includegraphics[width=\columnwidth]{naca_tn_3650_fin_wave_drag}
\caption{Fin wave drag for a 60-deg delta wing: present DATCOM~4.1.5.1
implementation versus NACA~TN~3650~\cite{ulmann1956} free-flight data.
Mean absolute percentage error $= 21.0\%$ over 12~points, $M = 1.1$--$2.5$.}
\label{fig:fin-wave}
\end{figure}

\subsection{Base Drag}

Four submodels cover the base-drag regimes of a slender finned rocket. For
$M > 1.3$ with a turbulent boundary layer the supersonic base-drag correlation
%
\begin{equation}
C_{d,\,\mathrm{base}} = 0.064 + \frac{0.186}{M^2}
\label{eq:base-turb}
\end{equation}
%
is applied. This empirical form is the supersonic asymptote for slender
axisymmetric afterbodies documented in ESDU~77021~\cite{esdu77021}, and it was
validated here against NACA~TN~3393~\cite{chapman1955} at four turbulent points
over $M = 2.73$--$4.48$, giving a mean absolute percentage error of $15.9\%$
(Fig.~\ref{fig:base-pressure}). For perfect-finish vehicles with delayed
transition the laminar base-pressure correlation of Chapman~\cite{chapman1950}
%
\begin{equation}
C_{p,\,b,\,\mathrm{lam}} = \frac{C_{\mathrm{lam}}}{M^2\sqrt{\mathrm{Re}_L}} ,
\qquad C_{\mathrm{lam}} = 1870 ,
\label{eq:base-lam}
\end{equation}
%
is used, validated against the four laminar TN~3393 points with a mean absolute
percentage error of $4.4\%$. Applying the turbulent form, Eq.~\eqref{eq:base-turb},
to the laminar data instead yields a mean absolute percentage error of $44\%$:
the boundary-layer state at the base materially governs the closure. A
Chapman--Korst free-shear-layer model with the ESDU~77021~\cite{esdu77021}
boundary-layer-thickness correction blends with Eq.~\eqref{eq:base-turb} over
$M = 1.2$--$1.4$. A Viswanath boattail correction~\cite{viswanath1996} reduces
base drag by $15$--$40\%$ for typical boattail half-angles of $6$--$16$~deg, and a
power-on multiplier derived from the solid-rocket motor exit-flow guidance of
NASA~SP-8039~\cite{nasa_sp8039} reduces base drag during burn using the nozzle
exit-area and pressure ratios. Below $M = 0.85$ a Hoerner subsonic
correlation~\cite{hoerner1965} is recovered, and the transonic peak near
$M \approx 1.05$ is represented by a degree-4 polynomial anchored at the regime
boundaries.

\begin{figure}[t]
\centering
\includegraphics[width=\columnwidth]{naca_tn_3393_base_pressure}
\caption{Turbulent and laminar base drag versus
NACA~TN~3393~\cite{chapman1955}. The supersonic turbulent correlation,
Eq.~\eqref{eq:base-turb}, attains a mean absolute percentage error of $15.9\%$;
the Chapman laminar correlation~\cite{chapman1950}, Eq.~\eqref{eq:base-lam},
attains $4.4\%$.}
\label{fig:base-pressure}
\end{figure}

\subsection{Skin Friction: Van Driest II}

The compressible skin-friction coefficient uses the Van~Driest~II transformation
of Eqs.~\eqref{eq:vd2_defs}--\eqref{eq:vd2_fc} (Sec.~\ref{sec:atmosphere}), in the
Hopkins--Inouye formulation~\cite{hopkins1971}, replacing the legacy Eckert
reference-temperature method. The transformed Reynolds number
$F_x\,\mathrm{Re}_x$, with $F_x = F_\theta/F_c$, is substituted into the
Karman--Schoenherr implicit relation
%
\begin{equation}
\frac{0.242}{\sqrt{C_{f,\,\mathrm{inc}}}}
   = \log_{10}\!\left(\mathrm{Re}_x\,C_{f,\,\mathrm{inc}}\right) ,
\label{eq:karman-schoenherr}
\end{equation}
%
solved by fixed-point iteration seeded with Schlichting's
$\log^{-2.58}$ approximation, and the compressible coefficient follows as
$C_f = C_{f,\,\mathrm{inc}}/F_c$. Hopkins and Inouye~\cite{hopkins1971}
demonstrated that Van~Driest~II gave the best agreement with experiment among
candidate transformations over $M = 1.5$--$9$. At $M = 5$ the compressible $C_f$
is approximately half of its incompressible value at matched length Reynolds
number (Fig.~\ref{fig:vd2}); the Eckert reference-temperature method reproduces
this trend to within roughly $10\%$ over $M = 1.5$--$9$, and Van~Driest~II was
retained for its slightly closer agreement with experiment at the upper end of
this range~\cite{hopkins1971}.

\begin{figure}[t]
\centering
\includegraphics[width=\columnwidth]{van_driest_ii_cf_vs_mach}
\caption{Van~Driest~II compressible skin-friction coefficient versus Mach number
for a representative slender body at length Reynolds number
$\mathrm{Re}_L = 10^7$, after the Hopkins--Inouye
formulation~\cite{hopkins1971}.}
\label{fig:vd2}
\end{figure}

\subsection{Boundary-Layer Transition}

The transition Reynolds number is Mach-dependent, with a laminar-fraction cap
that prevents unphysically long laminar runs at high Mach. The low-subsonic
crossover $\mathrm{Re}_{x,\mathrm{tr}} = 5\times10^5$ rises to approximately
$5\times10^6$ at $M = 4$, consistent with the compressible-stabilization trend
of the boundary layer. The laminar fraction is capped at $0.7$ of the body
length to preserve consistency with the post-shock entropy layer.

\subsection{Hypersonic Blending}

Above $M = 5$ the inviscid nose and body pressure distributions transition to
Modified Newtonian theory, $C_p = C_{p,\max}\sin^2\theta$, with $C_{p,\max}$
taken from the normal-shock stagnation pressure of Sec.~III. A smoothstep weight
$w = 3t^2 - 2t^3$ with $t = (M - 4)/2$ blends the shock-expansion result into the
Newtonian asymptote over $M = 4$--$6$. The combined hypersonic model was
assessed against DTIC~AD0487365~\cite{grabow1965} ballistic-range cone foredrag
at $M = 6.5$--$17.2$ for 8-, 12-, and 16-deg half-angle cones, giving an
aggregate mean absolute percentage error of $19.7\%$ across 11~points
(Fig.~\ref{fig:cone-drag}). The largest residual, $+57\%$, occurs for the
thinnest 8-deg cones at the lowest-Reynolds-number row, where friction and base
drag dominate the foredrag and the reference boundary-layer state is incompletely
specified; the blunter 16-deg cone agrees most closely. These results are
reported as exploratory: the headline flight validation (Sec.~VIII) extends only
to $M = 4.33$, and the hypersonic closure is a demonstrated method capability
rather than a flight-validated claim.

\begin{figure}[t]
\centering
\includegraphics[width=\columnwidth]{hypersonic_cone_drag}
\caption{Hypersonic cone foredrag versus DTIC~AD0487365~\cite{grabow1965}
ballistic-range data, $M = 6.5$--$17.2$, half-angles $8$, $12$, and $16$~deg.
Aggregate mean absolute percentage error $= 19.7\%$ over 11~points; the largest
residual is $+57\%$ for the thinnest cones at the lowest Reynolds number.}
\label{fig:cone-drag}
\end{figure}

\begin{table*}[t]
\centering
\caption{Drag submodel inventory (present work).}
\label{tab:drag-inventory}
\begin{tabularx}{\textwidth}{l l X l l}
\toprule
Submodel & Regime & Source / formulation & Validation source & Metric \\
\midrule
Taylor--Maccoll cone pressure & $M = 1$--$17$ & Exact conical-flow ODE (Sec.~III)
   & NACA~RM~A52H28 cone case & Exact (analytic) \\
Shock-expansion ogive & $M = 1.3$--$10$ & Eq.~\eqref{eq:strip-wave}, 100-strip integrator
   & NACA~RM~A52H28~\cite{a52h28} & MAE $0.029$ (5 shapes) \\
Dahlem--Buck shape factor & $M = 1.3$--$6$ & Eq.~\eqref{eq:dahlem-buck}
   & NACA~RM~A52H28~\cite{a52h28} & $K_{\mathrm{shape}}$ within $5\%$ \\
Fin wave drag, DATCOM 4.1.5.1 & $M = 0.9$--$5{+}$ & Eq.~\eqref{eq:datcom-fin}
   & NACA~TN~3650~\cite{ulmann1956} / Ackeret & MAPE $21\%$; analytic $0.00\%$ \\
Supersonic turbulent base & $M > 1.3$ & Eq.~\eqref{eq:base-turb}
   & NACA~TN~3393~\cite{chapman1955} (turb.) & MAPE $15.9\%$ \\
Chapman laminar base & $M = 1.3$--$4.5$ & Eq.~\eqref{eq:base-lam}
   & NACA~TN~3393~\cite{chapman1955} (lam.) & MAPE $4.4\%$ \\
Chapman--Korst shear layer & $M = 1.2$--$1.4$ & ESDU~77021~\cite{esdu77021} BL correction
   & Blend region & Continuity \\
Viswanath boattail & Any $M$, boattail & $\Delta C_{d,b}(\theta_{bt}, M)$
   & Viswanath~\cite{viswanath1996} & $15$--$40\%$ reduction \\
Power-on base & During burn & $k_{po}(A_e/A_{\mathrm{ref}}, p_e/p_\infty)$
   & NASA~SP-8039~\cite{nasa_sp8039} & Qualitative \\
Van~Driest~II $C_f$ & $M = 1.1$--$9$ & Eqs.~\eqref{eq:vd2_defs}--\eqref{eq:vd2_fc} (Sec.~\ref{sec:atmosphere}), \eqref{eq:karman-schoenherr}
   & Hopkins \& Inouye~\cite{hopkins1971} & ${\sim}50\%$ $C_f$ at $M = 5$ \\
Modified Newtonian & $M > 5$ & $C_p = C_{p,\max}\sin^2\theta$
   & DTIC~AD0487365~\cite{grabow1965} & MAPE $19.7\%$ \\
\bottomrule
\end{tabularx}
\end{table*}
